\documentclass{situationist-preprint}

\title{Ineffability and Structure}
\msid{MND-10-b}
\msfamily{MND-10}
\msdiscipline{Philosophy of Mind}
\mstype{Companion manuscript}
\mskeywords{qualia; ineffability; relational description; quality space; Newman's objection}

\begin{document}
\maketitle

\begin{abstract}
Qualia are said to be ineffable, but much about them can evidently be said: that red is more similar to orange than to green, that a sound is louder than another. We give a precise sense to ineffability using the result of MND-10 that a relational description of a quality space is complete up to structure-preserving permutation. Ineffability is the impossibility of conveying the character of a quality, as opposed to its structure, by description. We show that this is a consequence of the intrinsicness of qualia, that it is compatible with the full describability of their structure, and that it explains why qualia seem ineffable without making them mysterious.
\end{abstract}

\section{What can be said}

MND-10 argues that qualia exist, that they are intrinsic to the experiences that have them, and that the structure of each modality's qualia can be fully described in relational terms (MND-10-a, Theorem~5.4). A quality space is a set of qualities together with relations such as similarity, betweenness and order of intensity (Definition~3.3). Much of what ordinary talk says about qualities is relational: that one hue is between two others, that one pain is worse than another, that a taste is like another but sharper. All of this can be conveyed to someone who has never had the experiences in question. If ineffability meant that nothing could be said, qualia would not be ineffable, and the claim that they are has been held to show that the notion is confused \citep{Dennett1988}.

\begin{definition}\label{d:ineff}
A quality $q$ is \emph{ineffable} if no description, in terms available to someone who has not had $q$ or a quality of the same intrinsic character, conveys which intrinsic character $q$ has.
\end{definition}

The definition locates ineffability in the character of a quality and not in its structure. It concerns what a description can convey, not what can be thought: a subject who has had $q$ can think of it under a phenomenal concept, and can describe its relations to other qualities in public terms.

\section{Ineffability from intrinsicness}

\begin{proposition}\label{p:ineff}
Every quale is ineffable in the sense of Definition~\ref{d:ineff}, and every structural fact about a quality space can be conveyed by description.
\end{proposition}

\begin{proof}
A description available to someone who has not had the quality is a description in relational or physical terms; it does not use a phenomenal concept of the quality. By Proposition~6.1 of MND-10-a, a complete relational description of a quality space determines the assignment of qualia to places exactly up to its automorphisms, and never determines which intrinsic character the occupant of a given place has. By Lemma~5.1, a relational description is satisfied equally by any domain with the right structure, so no relational description distinguishes the actual intrinsic character of a quality from another that could occupy its place. Therefore no such description conveys which intrinsic character a quality has. The second claim is Theorem~5.4 of MND-10-a.
\end{proof}

The proposition separates two things the tradition ran together. What is ineffable is not the quality as a whole but its character; its place in the space of qualities is fully effable. The ineffability follows from the same source as the intrinsicness of qualia, Newman's point that relational descriptions fix only structure \citep{Newman1928}, and from nothing more.

The proposition also explains the phenomenon from which the doctrine of ineffability arose. When we try to say what red is like, we find ourselves saying what it resembles, what it contrasts with and what it is associated with. Each of these is a relation. We sense that something has been left out, and Proposition~\ref{p:ineff} identifies what: the character of the occupant of the place we have described. The sense of a residue is accurate.

\begin{objection}
If the character of a quality cannot be conveyed, it cannot make a difference to anything that can be said or done, and qualia are epiphenomenal.
\end{objection}

\begin{reply}
What cannot be conveyed by description is which character a quality has, not whether it makes a difference. A quality's character can cause the subject's recognition, judgment and report; what the report cannot do is convey the character to a hearer who has not had it. A colour patch causes a photograph of itself, and a description of the photograph in terms of relations among its parts does not convey which colours it contains. MND-10-a, Objection~4, gives the general reply.
\end{reply}

\begin{objection}
Two subjects with inverted qualities would describe their quality spaces identically. If ineffability is as stated, the difference between them is undetectable, and it is idle to posit it.
\end{objection}

\begin{reply}
The first claim is correct and is a consequence of Proposition~\ref{p:ineff}: an inversion that preserves relations is undetectable by any relational test, which is the conclusion of MND-22 \citep{Shoemaker1982}. The second does not follow. That a difference cannot be detected by description does not show that there is no difference; it shows only that the difference lies in what description cannot convey, which is what intrinsicness predicts.
\end{reply}

\section{Conclusion}

Qualia are ineffable in a precise sense: no description in public terms conveys the intrinsic character of a quality, although every structural fact about it can be conveyed. Ineffability is the shadow of intrinsicness, and the completeness of relational description up to automorphism marks its exact extent.

\begin{thebibliography}{9}\small
\bibitem[Dennett(1988)]{Dennett1988} Dennett, D.~C. (1988). Quining qualia. In A.~J. Marcel and E.~Bisiach (eds.), \emph{Consciousness in Contemporary Science}. Oxford University Press.
\bibitem[Newman(1928)]{Newman1928} Newman, M.~H.~A. (1928). Mr.\ Russell's \enquote{causal theory of perception}. \emph{Mind}, 37(146), 137--148.
\bibitem[Shoemaker(1982)]{Shoemaker1982} Shoemaker, S. (1982). The inverted spectrum. \emph{The Journal of Philosophy}, 79(7), 357--381.
\end{thebibliography}

\end{document}
